\documentclass[10pt,a4paper]{amsart}
\usepackage[margin=19mm]{geometry}
\usepackage{amsmath,amssymb,mathtools}
\usepackage[scr=rsfs,cal=euler]{mathalfa}
\usepackage{xfrac}
\usepackage[unicode,hidelinks,bookmarks=false]{hyperref}
\newcommand{\ie}{i.e.\ }
\newcommand{\cf}{cf.\ }
\newcommand{\resp}{respectively\ }
\newcommand{\Subsection}{\subsection}
\providecommand{\nobreakdash}{\nobreak\hbox{-}}
\setlength{\emergencystretch}{2em}
\multlinegap0pt
\usepackage{enumerate}
\usepackage{amssymb}
\usepackage{algorithm}\usepackage{algpseudocode}
\usepackage{tikz}
\newtheorem{prblm}{Problem}
\newtheorem{rmrk}{Rmrk}
\newtheorem{theorem}{Theorem}
\usepackage{cleveref}
\crefname{prblm}{Problem}{Problems}
\crefname{rmrk}{Rmrk}{Rmrks}
\crefname{theorem}{Theorem}{Theorems}
\newcommand{\Ap}{{\bt{A}}}
\newcommand{\Cp}{{\bt{C}}}
\newcommand{\Nr}{{N_r}}
\newcommand{\R}{\mathbb{R}}
\newcommand{\bt}[1]{{\mathbf{#1}}}
\newcommand{\bv}[1]{{\mathbf #1}} 
    \newcommand{\com}[1]{{{#1}_{\rm{c}}}} %complexity-reduced
\newcommand{\dreg}{{{\bv{d}}}}
\newcommand{\fopt}{{\eta}}
\newcommand{\fs}[1]{{{#1}}} %% function in function space
        \newcommand{\prol}[1]{{{#1}_{\rm{p}}}} %prolongated compressed
\newcommand{\rom}[1]{{{#1}_{\rm{r}}}}    
\newcommand{\tL}{{\fs{\varphi}}}
        \newcommand{\thi}[1]{{{#1}_{\rm{t}}}}	% thin compressed
\newcommand{\vAp}{\bv{a}}
\newcommand{\vwei}{\bv{u}} % vector quadrature weight
\newcommand{\wei}{{\omega}} % quadrature weight
\title{Reducing Training Complexity in Empirical Quadrature-Based Model Reduction via Structured Compression}
\author{Bj\"orn Liljegren-Sailer}
\date{Source: arXiv:2512.14416v2 (CC BY 4.0)}
\begin{document}
\maketitle

\noindent\emph{Source and licence.} Reducing Training Complexity in Empirical Quadrature-Based Model Reduction via Structured Compression. Version arXiv:2512.14416v2 (CC BY 4.0), distributed by arXiv under the Creative Commons Attribution 4.0 International licence, http://creativecommons.org/licenses/by/4.0/, as recorded in the arXiv licence metadata for this exact version and shown on the abstract page https://arxiv.org/abs/2512.14416v2. Full licence notice: \texttt{licenses/arxiv-2512.14416-CC-BY-4.0.txt}.\par\smallskip

\noindent\textit{Editorial scope.} Counted unit: the article's theorem the:CompressedOptFun (source0.tex lines 873--883). The article's complete original proof of it is delivered with it (source0.tex lines 884--908, one \texttt{proof} environment of 25 lines). No mathematics is written, repaired or re-derived: the statement and the proof are the article's own lines, in the article's own notation, and no retained line is substituted or re-flowed.  Retained uncounted as the source-local context the proof needs: lines 421--432; lines 488--535; lines 634--634; lines 859--861; lines 867--869.  Nothing was omitted from any retained span, so the package carries no omission record.  No retained line contains a citation, so the wrapper carries no bibliography and no bibliography entry.  No retained line contains a figure environment, an image-inclusion command or drawing code, so no image is delivered and the package declares no asset dependency.  Editorial formatting only, in addition: an AMS \texttt{amsart} preamble that restates the article's own declarations and macros used by the retained lines, namely the environments \texttt{prblm}, \texttt{rmrk}, \texttt{theorem} on separate counters, together with the standard packages those lines need.  The retained environments print in this document as Prblm 1, Rmrk 1, Theorem 1 in order of appearance, whereas the article prints the same items with the numbering of its own full document.  The complete native source of the article as distributed in the arXiv e-print for this exact version is bundled unchanged as \texttt{sources/191097/source0.tex}.
\begin{prblm}\label{prob:unc-CrTrain}
%%\text{ $ $}\\[-1.5cm]
\begin{align*}
\text{Solve} \quad & \min_{\substack{I_c \subset \{1,\ldots,M\},\, |I_c| \leq \com{M} \\ \vwei=[\wei_1,\ldots,\wei_M]^T}} 
\| \tilde{\Ap}(\vwei - \tilde{\vwei}) \|^2 + (\dreg^T(\vwei - \tilde{\vwei}))^2 \\
& \text{s.t.} \quad \wei_m \geq 0 \text{ for } m \in I_c,\quad \wei_{\bar{m}} = 0 \text{ for } \bar{m} \notin I_c.
\end{align*}
\end{prblm}
\begin{rmrk}\label{rem:standard-least-squares}
Note that the cost function of \Cref{prob:unc-CrTrain} can be rewritten in the standard least-square form $F: \vwei \mapsto \| \mathcal{A} \vwei -\bv{g} \|^2$, defining  $\mathcal{A} =  [\tilde{\Ap}^T, \dreg]^T$ and  $\bv{g} = \mathcal{A} \tilde{\vwei}$. 
\end{rmrk}
%%Thus, we have a non-negative least squares problem with a sparsity constraint, which we supplemented by a regularization term.

\subsection{The preprocessing problem}\label{subsec:preprocessing-ansatz}
%%
When $K$ is compressed to $\thi{K} \ll K$ in our approach, this yields a compressed solution manifold matrix $\thi{\Ap}$ and a respective compressed cost function of the form
%This relates to replacing the cost functional $\fopt(\vwei)$ in \Cref{prob:unc-CrTrain} by a compressed version, given by
\begin{align} \label{eq:foptThin}
	\thi{\fopt}(\vwei) = \| \thi{\Ap}( \vwei - \tilde{\vwei}) \|, \hspace{1cm} \text{with } \thi{\Ap}\in \R^{\thi{K}\Nr , M}.
\end{align}
The compressed training problem is then obtained replacing the exact cost function in \Cref{prob:unc-CrTrain} by its compressed counterpart.
%%
\begin{prblm} \label{prob:con-CrTrain}
%%Let the  compressed function $\thi{\fopt}$ be defined as above for a given matrix~$\prol{\Ap}$.
Let $\thi{\fopt}$ be as in \eqref{eq:foptThin} and let $\tilde{\vwei}$ denote the truth weights.
\begin{align*} 
\text{Solve \quad}	 &\min_{ \substack{I_c\subset \{1,\ldots M\}, \, |I_c | \leq \com{M} \\  \vwei=[\wei_1,\ldots,\wei_M]^T } }  \hspace{0.2cm} \thi{\fopt}(\vwei)^2 + (\dreg^T( \vwei -  \tilde{\vwei}))^2
 \\[-0.1cm]
	& \hspace{0.3cm} {\small \smash{ \text{s.t.}  \quad \wei_m \geq 0 \, \text{ for } m \in I_c, \quad \wei_{\bar{m}}  = 0 \text{ for } \bar{m}  \notin I_c  }},
\end{align*}
\end{prblm}
In what follows, we motivate and derive the compression step itself.
%%
We start by noting that $\tilde{\Ap} \in \mathbb{R}^{K\Nr , M}$ is obtained from three-dimensional data and represents one matricization of a tensor in $\mathbb{R}^{K , \Nr , M}$. For structured analysis, however, a different matricization is preferable, defined as
{\small
\begin{align}
\begin{split}\label{eq:C-data}
\tilde{\Cp} &=
\begin{pmatrix}
\beta_{\rm *}^1\!\left(f(\bv{x}^1), \rom{\tL}^{\!1}\right) & \ldots & \beta_{\rm *}^1\!\left(f(\bv{x}^K), \rom{\tL}^{\!1}\right) \\
\beta_{\rm *}^2\!\left(f(\bv{x}^1), \rom{\tL}^{\!1}\right) & \ldots & \beta_{\rm *}^2\!\left(f(\bv{x}^K), \rom{\tL}^{\!1}\right) \\
\vdots & \ddots & \vdots \\
\beta_{\rm *}^M\!\left(f(\bv{x}^1), \rom{\tL}^{\!1}\right) & \ldots & \beta_{\rm *}^M\!\left(f(\bv{x}^K), \rom{\tL}^{\!1}\right) \\
\vdots & \ddots & \vdots \\
\beta_{\rm *}^M\!\left(f(\bv{x}^1), \rom{\tL}^{\!\Nr}\right) & \ldots & \beta_{\rm *}^M\!\left(f(\bv{x}^K), \rom{\tL}^{\!\Nr}\right)
\end{pmatrix} \\
&=
\begin{pmatrix}
\vAp^{1,1} & \ldots & \vAp^{K,1} \\
\vdots & \ddots & \vdots \\
\vAp^{1,\Nr} & \ldots & \vAp^{K,\Nr}
\end{pmatrix}.
\end{split}
\end{align}
}%%
%%
By construction, $\tilde{\Cp}$ contains the same entries as $\tilde{\Ap}$. To fully exploit the underlying structure, we employ a factorization 
\begin{align*}
\tilde{\Cp} = {\bt{N}} \hat{\bt{G}}, \hspace{1cm} \bt{N} \in \R^{N_r M, M_J}, \quad \hat{\bt{G}} \in \R^{M_J,K}
\end{align*}
with $M_J \geq M$, and where $\bt{N}$ is a high-dimensional but sparse matrix encoding the constraints inherent in the data. As will be shown, the optimization can be reduced to a problem of the dimension of $\hat{\bt{G}}$, which is significantly lower than the dimension of $\tilde{\Ap}$. The derivation of the aforementioned factorization of $\tilde{\Cp}$ depends on the specific training problem; details are provided in Section~\ref{sec:separatedData}. Setting these details aside, the structured training problem of our proposed preprocessing can now be stated.

\section{Structured representation of the training data} \label{sec:separatedData}

\begin{align} \label{eq:costFunProl}
	\thi{\fopt}(\vwei) = \| \prol{\Ap}( \vwei - \tilde{\vwei}) \|, \hspace{1cm} \text{with } \prol{\Ap}\in \R^{{K}\Nr , M},
\end{align}

\begin{align} \label{eq:kap-error}
\kappa = \| \tilde{\Ap} - \prol{\Ap} \|_{\rm{F}} = \|  \tilde{\Cp} - \bt{N} {\bt{G}}_p^* \|_{\rm{F}} = \sqrt{\sum_{i > \thi{K}} \sigma_i^2}
\end{align}

\begin{theorem}\label{the:CompressedOptFun}
Consider \Cref{prob:unc-CrTrain} and \Cref{prob:con-CrTrain}. Let $\prol{\Ap}$ be defined as in~\eqref{eq:costFunProl}, and let $\kappa$ be as in~\eqref{eq:kap-error}. Define $d_{\rm min}$ as the smallest entry of $\dreg$, and assume $d_{\rm min} > 0$. Let $\vwei$ be a feasible point of \Cref{prob:con-CrTrain} such that $|\dreg^T(\vwei - \tilde{\vwei})| \leq \epsilon$ for some $\epsilon \geq 0$.

%%
Then the cost function $\fopt : \vwei \mapsto \|\tilde{\Ap}(\vwei - \tilde{\vwei})\|$ of \Cref{prob:unc-CrTrain} satisfies the following two bounds:
\begin{align*}
\fopt(\vwei)   &\leq  \thi{\fopt}(\vwei) + \kappa  ||\vwei - \tilde{\vwei}||,  &&  \text{(a posteriori bound)} \\
\fopt(\vwei)  &\leq \thi{\fopt}(\vwei) + \kappa  \left( \frac{\sqrt{\com{M}}}{d_{\rm{min}}} ( \epsilon+ \dreg^T \tilde{\vwei})  + \|\tilde{\vwei}\| ) \right).
&& \text{(a priori bound)}
\end{align*}
\end{theorem}

\begin{proof}
By definition of the cost functions, it follows
{\small
\begin{align*}
\fopt(\vwei) &=|| \tilde{\Ap} (\vwei - \tilde{\vwei}) || \leq || (\tilde{\Ap} - \prol{\Ap}) (\vwei - \tilde{\vwei}) || + || \prol{\Ap} (\vwei - \tilde{\vwei}) ||  \\
&=  || (\tilde{\Ap} - \prol{\Ap}) (\vwei - \tilde{\vwei}) || + \thi{\fopt}(\vwei) \leq || (\tilde{\Ap} - \prol{\Ap})|| \, ||\vwei - \tilde{\vwei}||  + \thi{\fopt}(\vwei) \\
&\leq  || (\tilde{\Ap} - \prol{\Ap})||_{\rm{F}} \, ||\vwei - \tilde{\vwei}||  + \thi{\fopt}(\vwei)  = \kappa ||\vwei - \tilde{\vwei}|| + \thi{\fopt}(\vwei),
\end{align*}
}%%
using the triangle inequality, sub-multiplicativity, and the consistency of the Frobenius norm with the spectral norm. This proves the a posteriori bound.

To proof the a priori bound, it remains to bound the term $||\vwei - \tilde{\vwei}||$ independently of $\vwei$. First note that
due to the non-negativity of $\vwei$ and $\dreg$, respectively, the one-norm of $\vwei$ fulfills
\begin{align*}
||\vwei||_1 = \sum_{m\in I_c} |\wei_m|  =  \frac{1}{d_{\rm{min}}} \left(\sum_{m\in I_c} d_{\rm{min}} \wei_m \right)  \leq  \frac{1}{d_{\rm{min}}}(\dreg^T \vwei ).
\end{align*}
Since $\vwei$ has at most $\com{M}$ nonzero entries, it can be identified with a vector in $\R^{\com{M}}$, implying the norm relation $||\vwei|| \leq \sqrt{\com{M}} ||\vwei||_1$. Using the triangle inequality, we obtain
\begin{align*}
\| \vwei - \tilde{\vwei} \| &\leq  \| \vwei \| + \|\tilde{\vwei}\|
\leq \sqrt{\com{M}}  \| \vwei \|_1  + \|\tilde{\vwei}\| 
 \leq      \frac{\sqrt{\com{M}}}{d_{\rm{min}}}(\dreg^T \vwei ) + \|\tilde{\vwei}\| 
\\ &\leq \frac{\sqrt{\com{M}}}{d_{\rm{min}}} ( \epsilon+ \dreg^T \tilde{\vwei})  + \|\tilde{\vwei}\| .
\end{align*}
The a priori bound now follows from inserting the latter bound on $\| \vwei - \tilde{\vwei} \|$ into the a posteriori bound.
\end{proof}

\end{document}
